% Source-content checks and remaining gates: paper/manuscript/source_manifest.json.
\section{Пов'язані роботи}\label{sec:related}

\paragraph{Лексичний, щільний і гібридний пошук.}
BM25, виведений з імовірнісної моделі релевантності, лишається стандартною функцією лексичного ранжування~\citep{robertson2009bm25}; archive-index використовує BM25, вбудований у SQLite FTS5~\citep{sqlitefts5}. Щільний пошук кодує запити й документи нейронними енкодерами і ранжує за подібністю векторів~\citep{reimers2019sbert,karpukhin2020dpr}. Бенчмарк BEIR показав, що BM25 --- стійкий базовий метод на різнорідних задачах без донавчання, а щільні моделі поза своєю предметною областю часто поступаються іншим підходам~\citep{thakur2021beir}. Гібридні системи поєднують обидва ранжування; злиття за оберненими рангами --- простий спосіб без навчання об'єднати ранжовані списки~\citep{cormack2009rrf}. Адаптивне правило, яке досліджено тут, не зливає оцінок чи рангів: воно обирає між режимом точної фрази і режимом «спершу зміст» за покриттям найкращого лексичного збігу (розділ~\ref{sec:system}).

\paragraph{Моделі вкладень.}
EmbeddingGemma --- відкрита текстова модель вкладень приблизно на 300 млн параметрів, побудована на Gemma~3 й оцінена на MTEB~\citep{vera2025embeddinggemma}; у статті використано її наступницю EmbeddingGemma~2 у випуску на Hugging Face, з промптами й вимогами до точності обчислень з картки моделі~\citep{google2026embeddinggemma2card}. Матрьошкове навчання представлень дає вкладення, префікси яких придатні й у меншій розмірності~\citep{kusupati2022matryoshka}; звідси 256-вимірний варіант. Multilingual E5 --- сімейство відкритих багатомовних текстових моделей вкладень, навчених контрастно приблизно на мільярді багатомовних текстових пар із подальшим донавчанням з учителем~\citep{wang2024me5}; його велика модель --- другий щільний базовий метод. Спільні простори вкладень для тексту і зображень поширилися завдяки CLIP~\citep{radford2021clip}; тут текст, зображення і мовлення ділять простір EmbeddingGemma~2.

\paragraph{Українська і російська мови.}
Українська і російська мають багату флексію; збіг лише поверхневих форм може пропускати документи з іншими формами терміна із запиту. pymorphy2 дає словниковий морфологічний аналіз для обох мов~\citep{korobov2015pymorphy}; archive-index використовує його підтримуване відгалуження pymorphy3~\citep{pymorphy3}. Публічні багатомовні бенчмарки пошуку, як-от MIRACL, охоплюють багато мов з оцінками релевантності від людей~\citep{zhang2023miracl}, натомість UA-PAR поєднує змішані формати документів, файли репозиторію, українську й російську та запити з описом документа, який людина пам'ятає. Це спеціально побудоване середовище архіву, а не репрезентативна вибірка особистих архівів. \uapar{} створено саме для цього (розділ~\ref{sec:uapar}).

\paragraph{Дані для мультимодального оцінювання.}
Crossmodal-3600 містить 3\,600 географічно різноманітних зображень із підписами людей 36 мовами, написаними без перекладу~\citep{thapliyal2022xm3600}; FLEURS --- n-сторонньо паралельне мовлення 102 мовами на основі речень FLoRes-101~\citep{conneau2023fleurs}. Їхні українські й російські пари «підпис--зображення» та «транскрипт--запис» задають цілі зіставлення з джерельних наборів. Створення підписів і записів людьми не є незалежним оцінюванням релевантності для нашого пошукового завдання.

\paragraph{Оцінювання з неповними і модельними оцінками.}
nDCG зі ступеневими виграшами введено в \citet{jarvelin2002ndcg}; ми обчислюємо її за допомогою ir-measures~\citep{macavaney2022irmeasures}. Пул глибини $k$ залишає документи поза пулом неоціненими; оцінювання на стиснутому списку, яке викидає з прогону неоцінені документи, --- один зі стандартних способів перевірити, наскільки це важить~\citep{sakai2007bpref}. Для значущості ми використовуємо парні рандомізаційні та бутстреп-процедури того типу, що порівнювали \citet{smucker2007significance}, а для узгодженості на порядковій шкалі --- лінійно зважену каппу~\citep{cohen1968kappa}. Великі мовні моделі вивчали як оцінювачів релевантності: є повідомлення, що вони добре передбачають уподобання користувачів пошуку~\citep{thomas2024llmjudge}, і є обговорення ризиків заміни людей-оцінювачів~\citep{faggioli2023perspectives}. У цій статті кілька модельних суддів від різних постачальників, модель, що писала запити, не входить до міток, і показано, як ранжування систем змінюється залежно від судді.

\paragraph{Керування особистою інформацією.}
Пошук на робочому столі за власними файлами людини досліджено в Stuff I've Seen, де об'єднано доступ до раніше бачених листів, файлів і вебсторінок~\citep{dumais2003sis}. archive-index має ту саму мету, але додає дослідницькі джерела (код, сесії агентів, журнали запусків), мультимодальний вміст, дві флективні мови і виміряне порівняння стратегій ранжування. Індекси наближеного пошуку найближчих сусідів, як-от HNSW~\citep{malkov2020hnsw} і бібліотека Faiss~\citep{douze2024faiss}, забезпечують наближений векторний пошук; archive-index шукає точно, і в розділі~\ref{sec:results} виміряно, де цей вибір перестає бути дешевим.
